\section{Introduction}\label{sec:intro}

Jagged competence combines strong performance with failures on particular
inputs.\footnote{\citet{gans2026} credits the term ``jagged intelligence''
to Andrej Karpathy; we use ``competence'' for performance in this task.}
Field experiments document an uneven AI frontier \citep{dellacqua2023}, and
\citet{gans2026} models how average quality can hide local unreliability. The
pattern is easy to observe and hard to explain: in a large system we rarely
know which intermediate quantities the task requires, or whether the system
computes, uses and keeps them, and without that knowledge ``jagged'' remains
a description.

We study the question in a task built from two known layers. Layer A computes
five per-slot sums from a board of records. Layer B uses the slot-1 sum's
category and a mode carried over from earlier boards to set the probabilities
of the next board's category. The same board can need different predictions
after different histories, so B cannot be computed from the current A alone;
in the terms of Six Birds Theory, B is a strict extension of A
(\cref{sec:strict}). A working hypothesis of that programme, motivated by its
picture of theory levels \citep[\S1.1]{tsiokos2026foundations} rather than
proved for trained networks, is that a network trained on an upper theory
comes to know the theory beneath it. Because both layers are known exactly,
we can test this: train networks on B only, and measure what they acquire of
A and how it relates to where B fails.

The answer, in this task, is that learning B gives the network a jagged
version of A, and that this explains part of the jagged competence in B. The
network makes readable the part of A that B needs while the rest becomes
less readable; in r12's uniform-sampling census, category misreads
concentrate near the slot-1 cutoffs; and when A is trained explicitly, it
comes out largely but not exactly right. The networks reach a natural KL below
$3\times10^{-4}$ bits, more than thirty times below our 0.01-bit bar, yet
fail on constructed histories (maximum law TV 0.25--0.26), and their category
misreads sit near the slot-1 cutoffs (\cref{sec:res-recover,sec:res-average}).
Here ``acquire'' means what we measure, readable by probes and approximately
computed when supervised; we do not demonstrate an internal layer, and the
jagged A explains part of B's failures, not all of them.

A second lesson is that knowing an intermediate quantity is not one thing.
A network can hold A in readable form without established exact computation,
compute it without delivering it to the prediction, deliver it in a form that
supports less reliable prediction, or lose it under further training; and
even an exact, delivered A leaves some failures in B. These are separate
achievements, not a required sequence, and matched comparisons isolate the
effects of connection, encoding and prediction gradients
(\cref{sec:res-noise,sec:res-access,sec:res-encoding,sec:res-erosion,sec:res-reliable}).

\paragraph{Contributions.}
\begin{itemize}
\item The finding that learning a strict extension gives a network only a
  jagged version of the lower theory, and that this explains part of
  the upper one's failure, with the unexplained cases stated
  (\cref{sec:results,sec:discussion}).
\item A decomposition of what it means for a network to know an intermediate
  quantity into five measurable achievements (recoverability, computation,
  access, preservation and reliable prediction), with matched experiments on
  access, encoding, preservation and target noise (\cref{sec:results}).
\item A layered task with exhaustive board ground truth and calibrated
  measurements (\cref{sec:task,sec:measurements}), and a provided ledger and
  code rebuilding every reported number, figure and table (\cref{app:compute}).
\end{itemize}
